\textbf{Решение домашнего задания (Теоретическая часть)}

\vspace{1em}
\textbf{1. Исходные данные для расчета теоретического спектра}

Из условия задания известны следующие расчетные параметры:
\begin{itemize}
    \item Отношение сигнал/шум: $a = 1$.
    \item Эффективное напряжение шума: $U_{\text{ш}} = 0.1 \text{ В}$.
    \item Комплексный параметр: $G_{U_{\text{ш}}0} K_{0\text{ БВЧ}}^2 = \frac{U_{\text{ш}}^2}{\Pi_{\text{ш}}}$.
\end{itemize}
Для численного расчета предположим, что шумовая полоса $\Pi_{\text{ш}} = 1 \text{ МГц} = 10^6 \text{ Гц}$ (стандартное значение из предыдущих пунктов вашей лабораторной работы).

\vspace{1em}
\textbf{2. Определение коэффициентов $b_1(a)$ и $b_2(a)$}

По графикам коэффициентов из раздаточного материала, для значения $a = 1$ кривые $b_1(a)$ и $b_2(a)$ пересекаются. Их приближенное значение в точке пересечения составляет:
\begin{equation*}
    b_1(1) \approx b_2(1) \approx 0.65
\end{equation*}

\vspace{1em}
\textbf{3. Расчет составляющих спектральной плотности огибающей}

Вычислим базовый множитель, общий для обеих составляющих (подставляя выражение, заданное в условии):
\begin{equation*}
    C = (4 - \pi) \frac{U_{\text{ш}}^2}{\Pi_{\text{ш}}} = (4 - 3.14159) \frac{0.1^2}{10^6} = 0.8584 \cdot \frac{0.01}{10^6} \approx 8.58 \cdot 10^{-9} \text{ В}^2/\text{Гц}
\end{equation*}

Тогда максимальная амплитуда обеих спектральных составляющих составит:
\begin{equation*}
    G_{\text{max}} = b_{1,2}(1) \cdot C \approx 0.65 \cdot 8.58 \cdot 10^{-9} \approx 5.58 \cdot 10^{-9} \text{ В}^2/\text{Гц}
\end{equation*}

Запишем выражения для прямоугольной и треугольной составляющих:
\textbf{Прямоугольная составляющая} ($0 \le f \le \Pi_{\text{ш}}/2$):
\begin{equation*}
    G_1(f) \approx 5.58 \cdot 10^{-9} \text{ В}^2/\text{Гц}
\end{equation*}
\textbf{Треугольная составляющая} ($0 \le f \le \Pi_{\text{ш}}$):
\begin{equation*}
    G_2(f) \approx 5.58 \cdot 10^{-9} \cdot \left(1 - \frac{f}{\Pi_{\text{ш}}}\right) \text{ В}^2/\text{Гц}
\end{equation*}

\vspace{1em}
\textbf{4. Суммарная спектральная плотность (Теоретическая)}

Суммарная теоретическая спектральная плотность огибающей $G_{\Sigma}(f) = G_1(f) + G_2(f)$ представляет собой кусочно-линейную функцию:
\begin{itemize}
    \item При $f = 0$: 
    $G_{\Sigma}(0) = G_1(0) + G_2(0) = 5.58 + 5.58 = 11.16 \cdot 10^{-9} \text{ В}^2/\text{Гц}$.
    \item При $f = \Pi_{\text{ш}}/2 = 0.5 \text{ МГц}$:
    \begin{align*}
        G_{\Sigma}(0.5-0) &= 5.58 + 5.58 \cdot 0.5 = 8.37 \cdot 10^{-9} \text{ В}^2/\text{Гц} \\
        G_{\Sigma}(0.5+0) &= 0 + 5.58 \cdot 0.5 = 2.79 \cdot 10^{-9} \text{ В}^2/\text{Гц} \quad \text{(возникает резкая ступенька)}
    \end{align*}
    \item При $f = \Pi_{\text{ш}} = 1 \text{ МГц}$: 
    $G_{\Sigma}(1) = 0$.
\end{itemize}

\begin{figure}[h]
    \centering
    \begin{tikzpicture}[scale=1.2]
        % Оси
        \draw[->, thick] (0,0) -- (6,0) node[right] {$f$, МГц};
        \draw[->, thick] (0,0) -- (0,6) node[above] {$G_{\Sigma}(f), 10^{-9} \text{ В}^2/\text{Гц}$};
        
        % Сетка и подписи осей
        \draw[step=1cm, gray, very thin, dashed] (0,0) grid (5.5, 5.5);
        \draw (2.5, 1pt) -- (2.5, -1pt) node[below] {$0.5$};
        \draw (5, 1pt) -- (5, -1pt) node[below] {$1.0$};
        \foreach \y in {2, 4, 6, 8, 10, 12} \draw (1pt, \y/2.15) -- (-1pt, \y/2.15) node[left] {\y};
        \node[below left] at (0,0) {0};
        
        % Суммарная спектральная плотность (Теория)
        \draw[blue, very thick] (0, 11.16/2.15) -- (2.5, 8.37/2.15);
        \draw[blue, very thick] (2.5, 8.37/2.15) -- (2.5, 2.79/2.15); % Ступенька
        \draw[blue, very thick] (2.5, 2.79/2.15) -- (5, 0) -- (5.5, 0);
        
        % Точки 
        \filldraw[blue] (0, 11.16/2.15) circle (1.5pt) node[above right, text=black] {$11.16$};
        \filldraw[blue] (2.5, 8.37/2.15) circle (1.5pt) node[above right, text=black] {$8.37$};
        \filldraw[blue] (2.5, 2.79/2.15) circle (1.5pt) node[above right, text=black] {$2.79$};
        
        \node[blue, right] at (3.5, 3.5) {Теоретическая огибающая};
        
    \end{tikzpicture}
    \caption{Теоретическая спектральная плотность огибающей при $a=1$}
\end{figure}

\textit{Вывод:} Полученный теоретический график со ступенчатым изломом на частоте $\Pi_{\text{ш}}/2$ следует визуально сопоставить с экспериментально снятой в модели спектрограммой, чтобы убедиться в адекватности используемой математической модели детектора.